% Part: applied-modal-logics
% Chapter: epistemic-logic
% Section: relational-models

\documentclass[../../../include/open-logic-section]{subfiles}

\begin{document}

\olfileid{aml}{el}{rel}

\olsection{సంబంధ నమూనాలు}

జ్ఞానసంబంధ తర్కాల ప్రాథమిక అర్థవిజ్ఞాన భావన సాధారణ
మోడల్ తర్కాల్లోని భావనలాంటిదే. సంబంధ నమూనాల్లో
లోకాల సమితి, ఏ \tetoken{ప్రతిజ్ఞావాక్య చరాలు}{!!{propositional
variable}s} ఏ లోకాల వద్ద ``సత్యం''గా పరిగణించాలో
నిర్ణయించే కేటాయింపు ఉంటాయి. ఒక్క కర్తను మాత్రమే
పరిగణిస్తే, మామూలుగా ఒకే ప్రాప్యత సంబంధం ఉంటుంది.
అయితే బహు-కర్త జ్ఞానసంబంధ తర్కంలో ఆ ఒక్క ప్రాప్యత
సంబంధం స్థానంలో ప్రాప్యత సంబంధాల సమితి ఉంటుంది;
మన కర్త-సంకేతాల సమితి~$G$లోని ప్రతి~$a$కు ఒకటి.

ఒక \emph{సంబంధ నమూనా}లో లోకాల సమితి ఉంటుంది;
ప్రతి కర్తకు ఒకటి చొప్పున ద్విస్థానిక ప్రాప్యత
సంబంధాలు ఆ లోకాలను కలుపుతాయి. ఏ
\tetoken{ప్రతిజ్ఞావాక్య చరాలు}{!!{propositional variable}s}
ఏ లోకాల వద్ద సత్యమో నిర్ణయించే కేటాయింపూ ఉంటుంది.

\begin{defn}
  బహు-కర్త జ్ఞానసంబంధ భాషకు ఒక \emph{నమూనా}
  $\mModel{M} = \tuple{W, R, V}$ అనే త్రయం; ఇందులో
  \begin{enumerate}
  \item $W$ అనేది ``లోకాల'' శూన్యం కాని సమితి,
  \item ప్రతి $a \in G$కు, ${R}_a$ అనేది~$W$పై
  ద్విస్థానిక ప్రాప్యత సంబంధం, మరియు
  \item $V$ అనేది ప్రతి \tetoken{ప్రతిజ్ఞావాక్య
    చరానికి}{!!{propositional
    variable}}~$p$ సాధ్య లోకాల సమితి $V(p)$ను
    కేటాయించే ప్రమేయకం.
  \end{enumerate}
  $R_a ww'$ అయితే, $w'$ను~$w$ నుంచి
  \emph{ఆ కర్తకు ప్రాప్యమైన} లోకం అంటాం.
  $w \in V(p)$ అయితే, $p$ అనేది~$w$ వద్ద
  \emph{సత్యం} అంటాం.
\end{defn}

ఇది సాధారణ మోడల్ తర్కంలోని విధానమే; అదనంగా
మరిన్ని ప్రాప్యత సంబంధాలు చేరాయి. ఇచ్చిన కర్త
విషయంలో, వారి ప్రాప్యత సంబంధం సాధారణంగా వారి
సమాచార స్థితిని సూచిస్తుందని అర్థం చేసుకుంటాం.
ఉదాహరణకు, $R_a ww'$ అంటే~$w$ వద్ద~$a$కు
ఉన్న సమాచారంతో~$w'$ అనుకూలంగా ఉందని తరచూ
తీసుకుంటాం. మరోలా చెప్పాలంటే,~$w$ వద్ద ఆ కర్తకు
లోకం~$w$కూ లోకం~$w'$కూ మధ్య తేడా గుర్తించలేరు.

\end{document}
